\section*{RESULTS}

\subsection{Circuit architecture: E/I modules as the oscillatory core}

How does a network of non-bursting neurons generate robust rhythmic output?
The \textit{Dendronotus} swim CPG achieves this through network hysteresis: slow synaptic feedback within contralateral E/I modules creates the history-dependence that individual neurons lack, driving alternation between left and right activity states (Fig.~\ref{fig:ei_mechanism}).
We constructed a computational model of this four-cell network by integrating validated two-cell subnetworks---half-center oscillators (HCOs) and E/I modules---following the approach developed in \cite{scully2024pairing}.

The circuit distributes rhythmogenic function across contralateral E/I modules rather than concentrating it in pacemaker cells (Fig.~\ref{fig:ei_mechanism}A).
Each module pairs an excitatory Si3 with a contralateral Si2 in a ``twisted'' configuration: Si3R drives Si2L and Si3L drives Si2R.
Within each module, Si3 provides excitatory drive through a logistic synapse with relatively fast kinetics, while Si2 returns slower inhibitory feedback through a logistic synapse---a synapse model with sigmoidal accumulation dynamics that captures the gradual buildup and delayed response observed in CPG synapses (see STAR Methods).
This delayed negative feedback loop constitutes the fundamental building block of the network.

Critically, each E/I module is ``almost-oscillatory'': isolated modules exhibit damped oscillations or subthreshold rhythmic fluctuations rather than sustained limit cycles.
These subcircuits lie close to oscillatory regimes in parameter space, near bifurcations that separate quiescent from rhythmic states, such that modest network coupling can recruit them into coordinated bursting.
This design principle contrasts with CPG architectures built from intrinsically bursting pacemaker neurons that generate robust oscillations independently of network context---here, no single component can generate rhythm autonomously.
The two HCOs (Si3L$\leftrightarrow$Si3R and Si2L$\leftrightarrow$Si2R) enforce anti-phase alternation but function primarily as stabilizers---the E/I modules supply the oscillatory drive.
Weak electrical coupling between E/I partners (Si3R$\leftrightarrow$Si2L and Si3L$\leftrightarrow$Si2R) provides rapid voltage coordination within each module and contributes to resilience.

The kinetic asymmetry between fast excitation and slow inhibition sets the burst timescale: Si3 excitation recruits Si2 quickly, but the inhibition that terminates Si3 firing accumulates over hundreds of milliseconds (Fig.~\ref{fig:ei_mechanism}B).
These timescales are not arbitrary modeling choices but reflect dynamics observed in the biological circuit.
Current injection into Si2R evokes dual-timescale responses in network partners: fast IPSPs in Si2L via the $\alpha$-synapse, immediate electrical deflections in Si3L, and slow accumulating inhibition from the logistic synapse that builds across hundreds of milliseconds (Fig.~\ref{fig:synaptic_timescales}).
Together, these dynamics establish the temporal framework for network bursting.
